\documentclass[twoside]{article}
\usepackage{package}
\usepackage[T1]{fontenc}
\usepackage{subfiles}
% Silenciar la advertencia de versalita + negrita inexistente
\WarningFilter{latex}{Some font shapes were not available}
\WarningFilter{latex}{Font shape `T1/lmr/bx/sc'}

% Silenciar la advertencia de caracteres perdidos en el nullfont (para 01_introduccion.tex)
\WarningFilter{latex}{Missing character: There is no}
%Nuevos comandos
\renewcommand{\thefootnote}{\arabic{footnote}}
\raggedbottom



\begin{document}
\begin{center}
    \thispagestyle{empty}
    \large{\textsc{Interfaces Gráficas y Especialidades en MATLAB}}\\[0.5cm] 
    \large{Fabrizio Benavidez\footnote[1]{\texttt{fabrizio.benavidez@mi.unc.edu.ar}}}\\[0.5cm]

     \href{https://github.com/Fabopy/Matlab-curso-}{%
        {\Large\faGithub}\quad
        \nolinkurl{https://github.com/Fabopy/Matlab-curso-}%
    }\\[0.5cm]

    \large{\footnotemark[1]\textit{Facultad de Ciencias Exactas, Físicas y Naturales, Ingeniería Aeronáutica}}
    \\[0.5cm]

    \begin{minipage}{0.7\textwidth}
    \hrule
    \vspace{0.2cm}
        \normalsize{\textsc{Resumen} En el siguiente trabajo se ``desarrolla"\ una aplicación hecha en MATLAB para ingresar numerador y denominador y poder seleccionar el tipo de entrada. En la misma aplicación te muestra las gráficas de las entradas que el usuario ingresa. Se utilizó la IA para desarrollar la guía 23 ya que el trabajo está basado en usar \textit{guide}, función que no existe más en MATLAB y pasó a llamarse de otra forma. 
        
        Por último, también se vieron algunos trucos o especialidades que tiene MATLAB: generador de números randoms sin que se repita alguno, diferencia entre predefinir el tamaño de un arreglo y cuándo no, tiempos de ejecución del código, definir el tamaño de una figura, cargar imágenes que tiene incorporado MATLAB, algunas funciones como el \texttt{why}, \texttt{bench} y \texttt{life}.}\\[0.2cm]
        \normalsize{\textit{Palabras claves}: MATLAB, aplicación, trucos\ }
    \vspace{0.2cm}
    \hrule
    \end{minipage}
    
    \vfil

    \includegraphics[width = 0.4\textwidth]{imagenes/Matlab_Logo.png}\\[0.5cm]

    \vfil

    \today\\[0.2cm]
    \large{Semana 11}
\end{center}







\newpage

\section*{Guía 23: Interfaces Gráficas en MATLAB}
\subfile{secciones/guia23}


\newpage

\section*{Guía 24: Especialidades en MATLAB} 
\subfile{secciones/guia24}














\newpage
% \thispagestyle{empty}
% \bibliography{biblioteca.bib}
% \bibliographystyle{plain}


\end{document}